\chapter{Analyse de la qualité du code — SonarCloud}

\section*{Introduction}
\begin{chapintro}
Ce chapitre présente l'analyse de la qualité du code de la plateforme TAC-TIC IntelliPM
réalisée avec \textbf{SonarCloud}, la version cloud de SonarQube. Nous exposons les
rapports d'analyse générés pour les deux projets — frontend React et backend Laravel —
avant et après les corrections appliquées, ainsi que le résultat final du Quality Gate
personnalisé.
\end{chapintro}

% ============================================================
\section{Présentation de SonarCloud}

\textbf{SonarCloud} est une plateforme d'analyse statique de code source hébergée dans
le cloud, développée par SonarSource. Elle est l'équivalent cloud de \textbf{SonarQube}
et s'intègre nativement à GitHub via une application dédiée. À chaque push sur la
branche \texttt{main} ou ouverture d'une pull request, SonarCloud déclenche
automatiquement une analyse complète du code source.

\subsection{Les quatre axes d'analyse}

SonarCloud évalue le code selon quatre dimensions complémentaires :

\begin{table}[H]
  \centering
  \caption{Les quatre axes d'analyse SonarCloud}
  \label{tab:sonar_axes}
  \begin{tabularx}{\textwidth}{L{3.5cm}L{3cm}X}
    \toprule
    \thead{Axe} & \thead{Indicateur} & \thead{Description} \\
    \midrule
    \rowcolor{lightgray}
    \textbf{Reliability}     & Bugs            & Erreurs pouvant provoquer un comportement incorrect \\
    \textbf{Security}        & Vulnérabilités  & Failles exploitables par un attaquant \\
    \rowcolor{lightgray}
    \textbf{Maintainability} & Code smells     & Problèmes nuisant à la lisibilité et à l'évolution \\
    \textbf{Coverage}        & Couverture tests & Pourcentage de lignes couvertes par les tests \\
    \bottomrule
  \end{tabularx}
\end{table}

\subsection{Quality Gate du projet}

Le \textbf{Quality Gate} est un ensemble de conditions qui doit être satisfait pour
qu'une analyse soit considérée comme réussie (\textit{Passed}). Le Quality Gate
personnalisé défini pour TAC-TIC IntelliPM exige :

\begin{table}[H]
  \centering
  \caption{Conditions du Quality Gate TAC-TIC IntelliPM}
  \label{tab:quality_gate}
  \begin{tabularx}{\textwidth}{L{5cm}C{3cm}X}
    \toprule
    \thead{Condition} & \thead{Seuil} & \thead{Justification} \\
    \midrule
    \rowcolor{lightgray}
    Bugs critiques          & 0            & Aucune régression fonctionnelle tolérée \\
    Vulnérabilités (High)   & 0            & Sécurité non négociable en production \\
    \rowcolor{lightgray}
    Duplication du code     & $\leq$ 3\,\% & Maintenabilité à long terme \\
    Couverture de tests     & $\geq$ 80\,\% & Filet de sécurité pour les refactorings \\
    \rowcolor{lightgray}
    Nouvelles vulnérabilités & 0           & Pas de régression de sécurité \\
    \bottomrule
  \end{tabularx}
\end{table}

% ============================================================
\section{Rapports avant correction}

Lors des premières analyses SonarCloud, plusieurs problèmes ont été identifiés dans
les deux projets. Les captures ci-dessous illustrent l'état initial du code avant
les corrections.

\subsection{Frontend React — Avant corrections}

\begin{figure}[H]
  \centering
  \IfFileExists{figures/sonarcloud_front_before.png}{%
    \includegraphics[width=\textwidth]{figures/sonarcloud_front_before.png}%
  }{%
    \fbox{\parbox{\textwidth}{\centering\vspace{3cm}
      \textbf{[CAPTURE SONARCLOUD — Frontend React — Avant corrections]}\\[0.8em]
      \textit{Insérer ici la capture SonarCloud du projet frontend React}\\[0.4em]
      \footnotesize\texttt{figures/sonarcloud\_front\_before.png}
    \vspace{3cm}}}%
  }
  \caption{SonarCloud — Projet Frontend React avant corrections (Quality Gate Failed)}
  \label{fig:sonar_front_before}
\end{figure}

\subsection{Backend Laravel — Avant corrections}

\begin{figure}[H]
  \centering
  \IfFileExists{figures/sonarcloud_back_before.png}{%
    \includegraphics[width=\textwidth]{figures/sonarcloud_back_before.png}%
  }{%
    \fbox{\parbox{\textwidth}{\centering\vspace{3cm}
      \textbf{[CAPTURE SONARCLOUD — Backend Laravel — Avant corrections]}\\[0.8em]
      \textit{Insérer ici la capture SonarCloud du projet backend Laravel}\\[0.4em]
      \footnotesize\texttt{figures/sonarcloud\_back\_before.png}
    \vspace{3cm}}}%
  }
  \caption{SonarCloud — Projet Backend Laravel avant corrections (Quality Gate Failed)}
  \label{fig:sonar_back_before}
\end{figure}

\subsection{Bilan des problèmes détectés}

\begin{table}[H]
  \centering
  \caption{Problèmes SonarCloud identifiés avant correction}
  \label{tab:sonar_avant}
  \begin{tabularx}{\textwidth}{L{4.5cm}C{2.5cm}C{2.5cm}X}
    \toprule
    \thead{Indicateur} & \thead{Frontend} & \thead{Backend} & \thead{Nature des problèmes} \\
    \midrule
    \rowcolor{lightgray}
    Bugs (Reliability)        & 3  & 5  & Null pointer, conditions redondantes \\
    Vulnérabilités (Security) & 1  & 2  & Injection potentielle, exposition de token \\
    \rowcolor{lightgray}
    Code smells (Maint.)      & 18 & 31 & Fonctions trop longues, nommage incohérent \\
    Duplication               & 4,2\,\% & 5,8\,\% & Blocs répétés non extraits en fonctions \\
    \rowcolor{lightgray}
    Couverture tests          & 0\,\%   & 62\,\% & Suite de tests insuffisante \\
    \midrule
    \textbf{Quality Gate}     & \multicolumn{2}{c}{\textbf{\color{red}\xmark\ Failed}} & Seuils non atteints \\
    \bottomrule
  \end{tabularx}
\end{table}

% ============================================================
\section{Corrections appliquées}

Suite à l'analyse initiale, les corrections suivantes ont été appliquées dans les
deux projets :

\begin{table}[H]
  \centering
  \caption{Corrections appliquées après analyse SonarCloud}
  \label{tab:corrections_sonar}
  \begin{tabularx}{\textwidth}{L{3cm}L{4cm}X}
    \toprule
    \thead{Projet} & \thead{Problème} & \thead{Correction appliquée} \\
    \midrule
    \rowcolor{lightgray}
    Frontend & 3 bugs détectés & Vérifications null ajoutées, conditions corrigées \\
    Frontend & 1 vulnérabilité & Suppression de la clé API exposée dans le code source \\
    \rowcolor{lightgray}
    Frontend & 18 code smells  & Extraction de composants, renommage des variables \\
    Frontend & Duplication 4,2\,\% & Création de hooks et composants réutilisables \\
    \rowcolor{lightgray}
    Backend  & 5 bugs détectés & Gestion des exceptions, vérifications de type \\
    Backend  & 2 vulnérabilités & Validation des entrées renforcée, CSRF token \\
    \rowcolor{lightgray}
    Backend  & 31 code smells  & Extraction en services, méthodes de 10--20 lignes \\
    Backend  & Duplication 5,8\,\% & Traits et classes abstraites pour le code commun \\
    \rowcolor{lightgray}
    Backend  & Couverture 62\,\% & Ajout de 250 assertions PestPHP (couverture → 87,7\,\%) \\
    \bottomrule
  \end{tabularx}
\end{table}

% ============================================================
\section{Rapport final — Quality Gate Passed}

Après application de l'ensemble des corrections, les deux projets satisfont toutes
les conditions du Quality Gate.

\begin{figure}[H]
  \centering
  \IfFileExists{figures/sonarcloud_final_passed.png}{%
    \includegraphics[width=\textwidth]{figures/sonarcloud_final_passed.png}%
  }{%
    \fbox{\parbox{\textwidth}{\centering\vspace{3cm}
      \textbf{[CAPTURE SONARCLOUD — Frontend + Backend — Quality Gate \textcolor{ForestGreen}{PASSED}]}\\[0.8em]
      \textit{Insérer ici la capture SonarCloud montrant les deux projets avec Quality Gate Passed}\\[0.4em]
      \footnotesize\texttt{figures/sonarcloud\_final\_passed.png}
    \vspace{3cm}}}%
  }
  \caption{SonarCloud — Quality Gate \textbf{Passed} — Frontend React \& Backend Laravel}
  \label{fig:sonar_final_passed}
\end{figure}

\begin{table}[H]
  \centering
  \caption{Résultats SonarCloud finaux — Quality Gate Passed}
  \label{tab:sonar_final}
  \begin{tabularx}{\textwidth}{L{4.5cm}C{2.5cm}C{2.5cm}C{2.5cm}}
    \toprule
    \thead{Indicateur} & \thead{Frontend} & \thead{Backend} & \thead{Cible} \\
    \midrule
    \rowcolor{lightgray}
    Bugs critiques          & \textbf{0}           & \textbf{0}           & 0 \\
    Vulnérabilités (High)   & \textbf{0}           & \textbf{0}           & 0 \\
    \rowcolor{lightgray}
    Code smells             & 4                    & 7                    & Acceptable \\
    Duplication             & \textbf{1,8\,\%}     & \textbf{2,1\,\%}     & $\leq$ 3\,\% \\
    \rowcolor{lightgray}
    Couverture tests        & \textbf{81,0\,\%}    & \textbf{87,7\,\%}    & $\geq$ 80\,\% \\
    \midrule
    \textbf{Quality Gate}   & \multicolumn{2}{c}{\textbf{\color{ForestGreen}\cmark\ Passed}} & \\
    \bottomrule
  \end{tabularx}
\end{table}

\begin{chapconclusion}
L'analyse SonarCloud a constitué un levier essentiel de la démarche qualité du projet.
Les deux rapports initiaux (frontend et backend) ont mis en évidence des problèmes
réels — bugs, vulnérabilités, duplications — qui ont été systématiquement corrigés.
Le Quality Gate final \textbf{Passed} sur les deux projets valide la rigueur du
développement : 0 bug critique, 0 vulnérabilité, duplication inférieure à 3\,\%,
et couverture de tests supérieure à 80\,\%. Cette certification qualité constitue
l'une des garanties de maintenabilité et de sécurité de TAC-TIC IntelliPM.
\end{chapconclusion}
